\section{Related Work}
\paragraph{Text Diffusion Models}
Early explorations of text diffusion models were based on a continuous space~\citep{li2022diffusion,gong2022diffuseq,DBLP:conf/iclr/ChenZH23}. Subsequently, discrete diffusion models~\citep{hoogeboom2021argmax,austin2021structured} directly introduced discrete noise to accommodate the discrete nature of text, demonstrating significant potential~\citep{Zheng2023ARD,lou2023discrete} and were further developed into mask diffusion models~\citep{shi2024simplified, ou2024your, sahoo2024simple}.
Recent work has explored scaling these models significantly, with DiffuLLaMA~\citep{gong2025scaling} being adapted from pretrained AR LLMs, and LLaDA~\citep{nie2024scaling} and Dream~\citep{dream2025} being the first open-source diffusion LLMs to achieve performance comparable to AR LLMs. Block diffusion~\citep{arriola2025block} proposes a hybrid approach that applies diffusion within each block~\citep{han-etal-2023-ssd}, serving as a midpoint between autoregressive and diffusion models. Multimodal models such as LaViDa~\citep{li2025lavida}, MMaDA~\citep{yang2025mmada}, and Dimple~\citep{yu2025dimple} combine text diffusion models with vision models. \citet{liu2025dllm, hu2025accelerating, ma2025dkv, fastdllm2025,sahoo2025esotericlanguagemodels} introduce caching and parallel decoding algorithms for dLLMs, significantly improving inference efficiency.

\paragraph{Code Generation}
Code generation is a crucial domain for LLMs~\citep{roziere2023code,sun2024survey}, exemplified by state-of-the-art open-source models like Qwen-2.5-Coder~\citep{hui2024qwen2} and OpenCoder~\citep{huang2024opencoder}, with wide applications in areas such as coding assistants and agents~\citep{xu2024lemur}.
CodeFusion~\citep{singh2023codefusion} was the first to combine diffusion models with code generation, but it was limited to small-scale models and simple tasks. Recent commercial-scale dLLMs, such as Mercury~\citep{inception2025mercury} and Gemini~\citep{geminidiffusion}, have demonstrated that diffusion-based code generators can achieve performance comparable to leading autoregressive code models while offering significantly faster generation speeds. 

\paragraph{Reinforcement Learning}
Reinforcement learning with verifiable reward (RLVR) using GRPO~\citep{shao2024deepseekmath,openr1,guo2025deepseek,bercovich2025llama} is highly effective in enhancing a language model's math reasoning~\citep{shao2025spurious} and code generation abilities~\citep{xie2025teaching}. \citet{wang2025octothinker} show the importance of mid-training during RL scaling. Combining RL and diffusion models, VRPO~\citep{li2025llada} introduces the efficient sampling algorithm from DPO~\citep{rafailov2023direct} for dLLMs.
d1~\citep{zhao2025d1} and MMaDA~\citep{yang2025mmada} optimize math reasoning in dLLMs using GRPO, but they rely heavily on block diffusion decoding during rollout and evaluation. LLadDou~\citep{huang2025reinforcing} trains an additional module to predict the rank score of tokens. For small text diffusion models, \citet{zhang2025target} propose the target concrete score matching framework, and \citet{zekri2025fine} introduces score entropy policy optimization (SEPO). Earlier, DDPO~\citep{black2023training} and DPPO~\citep{ren2024diffusion} formulated the diffusion process as a Markov Decision Process and performed policy optimization for continuous diffusion models.
